%% ZF+BPIT representation, with a separate countable ZF construction.

\section{Prime filters and Priestley duality}

\begin{annexitem}{annex:p10.priestley.1}{thm:p10-prime-filter-separation}{Prime-filter separation}
\begin{proof}
\paragraph{A maximal filter avoiding an ideal is prime.}
For a filter \(G\) and \(a\in D\), its extension generated by \(a\) is
\[
 G[a]=\{d:\exists g\in G\quad g\wedge a\leq d\}.
\]
The filter laws and leastness follow by taking meets of two
witnesses, just as for Boolean filters.
Suppose \(P\) is maximal among filters avoiding \(I\).
If \(a\vee b\in P\) but neither \(a\) nor \(b\) belongs to \(P\),
both generated extensions meet \(I\).
Thus there are \(p,q\in P\) and \(i,j\in I\) with
\[
 p\wedge a\leq i,\qquad q\wedge b\leq j.
\]
Distributivity gives
\[
 r=p\wedge q\wedge(a\vee b)
   =(p\wedge q\wedge a)\vee(p\wedge q\wedge b)
   \leq i\vee j.
\]
Here \(r\in P\) and \(r\in I\), a contradiction.
Also \(0\notin P\), because \(0\in I\).
Hence \(P\) is prime.
For finite \(D\), a maximal member of the nonempty finite set
of filters containing \(F\) and avoiding \(I\) exists in ZF.
This proves the finite case.

\paragraph{Finite generation and the propositional encoding.}
A bounded distributive sublattice generated by finitely many
elements \(d_1,\ldots,d_m\) is finite.
Every term in those generators is a join of meets of subsets
of the generators.
The empty meet and join represent \(1\) and \(0\).
The normal-form claim follows by induction on terms: join takes
the union of the finite collections of meets, while meet
distributes over both finite joins.
There are at most \(2^{2^m}\) such forms.

Use a classical proposition letter \(p_a\) for each \(a\in D\).
Let \(\Pi_D\) contain
\[
 p_1,\quad\neg p_0,\quad
 p_{a\wedge b}\leftrightarrow(p_a\wedge p_b),\quad
 p_{a\vee b}\leftrightarrow(p_a\vee p_b)
       \quad(a,b\in D).
\]
Add \(p_f\) for every \(f\in F\) and \(\neg p_i\) for every \(i\in I\).
Consider any finite subset of this classical theory.
Its mentioned lattice elements lie in a finite generated bounded
distributive sublattice \(D_0\).
The filter \(F\cap D_0\) and ideal \(I\cap D_0\) are disjoint.
The finite case supplies a prime filter of \(D_0\) extending the
former and avoiding the latter.
Its characteristic map satisfies the chosen clauses.
Set any unmentioned proposition letters to \(0\) to get a
valuation on the full alphabet.
Thus the theory is finitely satisfiable.

By arbitrary-variable propositional compactness under BPIT,
Theorem~\ref{thm:p02-arbitrary-prop-compactness-bpi}, it has a model \(v\).
The set \(P=\{a:v(p_a)=1\}\) contains \(1\), omits \(0\),
is meet-closed, and has the prime-join property.
If \(a\leq b\) and \(a\in P\), the meet equation for
\(a\wedge b=a\) forces \(b\in P\), so it is upward closed.
The added clauses give \(F\subseteq P\) and \(P\cap I=\varnothing\).
Only bounded-lattice equations were encoded, with no requirement
to preserve any additional Heyting implication.

\paragraph{The countable construction in ZF.}
For countable \(D\), fix an enumeration \(d_0,d_1,\ldots\),
with repetitions allowed, and let \(F_0=F\).
Recursively put
\[
 F_{n+1}=
 \begin{cases}
  F_n[d_n],&F_n[d_n]\cap I=\varnothing,\\
  F_n,&\text{otherwise}.
 \end{cases}
\]
At every stage the filter avoids the entire ideal \(I\).
The union \(P=\bigcup_n F_n\) is a filter: finitely many
arguments occur together at some stage.
It avoids \(I\), since each element of the union is in one stage.
If \(d_n\notin P\), the rejected extension at that stage meets
\(I\). Any filter containing \(P\) and \(d_n\) contains this
rejected extension, and therefore meets \(I\).
Thus \(P\) is maximal among filters containing \(F\) and avoiding
\(I\), and the first paragraph proves it prime.
Each recursion step is specified by a definite case distinction;
no sequence of arbitrary choices is needed.
Finite nonempty \(D\) admits an enumeration with repetitions.

Finally, if \(a\not\leq b\), the principal filter above \(a\)
and principal ideal below \(b\) are disjoint.
Apply the proved extension theorem to them.
For trivial \(D\), no disjoint filter and ideal exist, so the
general assertion remains valid vacuously.
\end{proof}
\end{annexitem}

\begin{annexitem}{annex:p10.priestley.2}{thm:p10-priestley-representation}{Priestley representation}
\begin{proof}
Prime-filter laws give
\[
 \widehat0=\varnothing,\quad\widehat1=X_D,\quad
 \widehat{a\wedge b}=\widehat a\cap\widehat b,\quad
 \widehat{a\vee b}=\widehat a\cup\widehat b.
\]
Hence the indicated differences form a basis.
Every \(\widehat a\) is a clopen upset.
If \(P\not\subseteq Q\), choose \(a\in P\setminus Q\);
then \(\widehat a\) separates them in the required order direction.
For distinct \(P,Q\), one of these noninclusions holds, so an
admissible clopen and its complement give disjoint neighborhoods.
The space is Hausdorff and totally order-disconnected.

\paragraph{Compactness.}
Let \(\mathcal U\) be an open cover without a finite subcover.
Take the set \(E\) of all pairs \((a,b)\) such that
\(\widehat a\setminus\widehat b\) is contained in a member of
\(\mathcal U\).
These basic sets cover \(X_D\).
No finite collection of them covers it, since choosing finitely
many parent cover members would then give a finite subcover.

Use the classical theory \(\Pi_D\) from the preceding annex
and add \(p_a\to p_b\) for every \((a,b)\in E\).
Every finite subset is satisfiable: choose a prime filter outside
the finitely many mentioned basic sets and use its characteristic
valuation. This satisfies all of \(\Pi_D\) as well as those
finitely many avoidance clauses.
BPIT and propositional compactness give a valuation of the entire
theory. Its prime filter belongs to none of the basic sets in
the covering refinement, a contradiction.
Thus \(X_D\) is compact.
The argument also shows directly why an empty space cannot
have an open cover without a finite subcover.

\paragraph{The map is an order embedding.}
The displayed identities make \(\eta_D\) a bounded lattice
homomorphism.
If \(a\leq b\), upward closure of filters gives
\(\widehat a\subseteq\widehat b\).
If \(a\not\leq b\), prime-filter separation gives a point in
\(\widehat a\setminus\widehat b\).
Thus order is reflected, and the map is injective.

\paragraph{All clopen upsets are represented.}
Let \(K\subseteq X_D\) be a clopen upset and fix \(P\in K\).
For each \(Q\notin K\), one has \(P\not\subseteq Q\), since
otherwise the upset property would put \(Q\) in \(K\).
The family
\(\{X_D\setminus\widehat a:a\in P\}\)
therefore covers \(X_D\setminus K\).
Together with \(K\), it covers the compact whole space.
A finite subcover supplies \(a_1,\ldots,a_n\in P\) such that
\[
 P\in\widehat{a_1\wedge\cdots\wedge a_n}\subseteq K.
\]
If no complements are needed, the empty meet is \(1\).
Consequently the collection of all \(\widehat a\subseteq K\)
covers \(K\); no simultaneous choice for each \(P\) is required.
Compactness, applied with the open complement of \(K\), reduces
this cover to a finite one.
Taking the corresponding finite join gives \(K=\widehat a\).
For \(K=\varnothing\), use \(a=0\).
This proves surjectivity.
In the trivial case \(X_D=\varnothing\), and both lattices have
one element, as asserted.
\end{proof}
\end{annexitem}

\begin{annexitem}{annex:p10.priestley.3}{thm:p10-priestley-reconstruction}{Recovering a Priestley space}
\begin{proof}
For \(x\in X\), membership at \(x\) is a bounded lattice
homomorphism from \(\operatorname{ClopUp}(X)\) to \(2\).
Its true elements form the prime filter \(\varepsilon_X(x)\).
If \(x\leq y\), every upset containing \(x\) contains \(y\),
so \(\varepsilon_X(x)\subseteq\varepsilon_X(y)\).
Conversely, if \(x\not\leq y\), order separation gives a
clopen upset in the former filter and absent from the latter.
Thus the map preserves and reflects order, and is injective.

Let \(P\) be a prime filter of the clopen-upset lattice.
The complementary set of lattice elements is an ideal:
it contains \(\varnothing\), is downward closed, and is join-closed
by primeness of \(P\).
Consider the closed family
\[
 P\ \cup\ \{X\setminus V:V\in\operatorname{ClopUp}(X),\ V\notin P\}.
\]
It has the finite-intersection property.
Indeed, an empty finite intersection would give
\(\bigcap_i U_i\subseteq\bigcup_j V_j\), with all \(U_i\in P\)
and all \(V_j\notin P\).
The left side belongs to \(P\), including the empty meet \(X\);
upward closure and primeness would then put one of the \(V_j\)
in \(P\), or put the empty join \(\varnothing\) there.
Both conclusions are impossible.
Compactness gives a point in the total intersection.
It lies in exactly the clopen upsets belonging to \(P\), so its
evaluation filter is exactly \(P\).
This proves surjectivity.

The preimages of the defining subbasic sets and their complements
are respectively the clopen upsets and their complements in \(X\).
Thus \(\varepsilon_X\) is continuous.
The prime-filter topology is Hausdorff by the elementary separation
argument at the start of Annex~\ref{annex:p10.priestley.2}, which
uses no BPIT.
It is a continuous bijection from a compact space to a Hausdorff
space and hence a homeomorphism.
For completeness, a compact subset of a Hausdorff space is closed:
separate an outside point from each point of the compact subset,
take all such separating neighborhoods, reduce the resulting
cover of the subset to a finite cover, and intersect the finitely
many neighborhoods of the outside point.
Thus a continuous image of a closed subset of \(X\) is closed,
proving continuity of the inverse.
Only finitely many choices enter that argument.
For empty \(X\), the clopen-upset lattice is trivial and has no
prime filters; the empty map is the required isomorphism.
\end{proof}
\end{annexitem}

\begin{annexitem}{annex:p10.priestley.4}{thm:p10-priestley-functors}{Functorial Priestley duality}
\begin{proof}
For a bounded lattice homomorphism \(h:D\to E\), the inverse
image of a prime filter contains \(1\), excludes \(0\), is
upward closed and meet-closed, and is prime by preservation of joins.
This proves \(X(h)\) is well-defined.
It is order-preserving under inclusion.
Moreover,
\[
 X(h)^{-1}[\widehat a]=\widehat{h(a)}.
\]
Preimages also commute with complements, so the equation proves
continuity on the generating subbasis.

For continuous order-preserving \(g:X\to Y\), the inverse image
of a clopen upset is a clopen upset.
The usual inverse-image laws preserve finite meets, joins,
and both bounds.
Both assignments preserve identities and reverse composition.

The two preceding representation theorems supply the isomorphisms
on objects.
For a lattice homomorphism \(h:D\to E\), the basic-preimage
equation gives
\[
 \operatorname{ClopUp}(X(h))\circ\eta_D=\eta_E\circ h.
\]
For \(g:X\to Y\), membership of a clopen upset \(U\subseteq Y\)
in the image of the point filter at \(x\) is equivalent to
\(x\in g^{-1}[U]\), hence to \(g(x)\in U\).
Therefore
\[
 X(\operatorname{ClopUp}(g))\circ\varepsilon_X
       =\varepsilon_Y\circ g.
\]
These are the naturality identities.

Finally, for \(P\in X_D\), pulling its point filter back along
\(\eta_D\) returns
\(\{a:P\in\widehat a\}=P\).
For a clopen upset \(U\subseteq X\), pulling \(\widehat U\)
back along \(\varepsilon_X\) returns
\(\{x:U\in\varepsilon_X(x)\}=U\).
The reconstruction composites are identities.
Thus the functors give a contravariant equivalence, including
the trivial lattice and empty ordered space.
\end{proof}
\end{annexitem}
